\Askhsh{36827}{2}
\begin{ylh}{1}
    \item 1.3 Μονότονες συναρτήσεις - Αντίστροφη συνάρτηση
    \item 2.2 Παραγωγίσιμες συναρτήσεις - Παράγωγος συνάρτηση
    \item 2.3 Κανόνες παραγώγισης
    \item 2.5 Θεώρημα Μέσης Τιμής Διαφορικού Λογισμού.
\end{ylh}
Θεωρούμε τις συναρτήσεις $f(x) = \ln x, x > 0$ και $g(x) = e^{-x}, x \in \mathbb{R}$.
\begin{alist}
    \item Να δικαιολογήσετε ότι η συνάρτηση $g$ έχει αντίστροφη και να αποδείξετε ότι $g^{-1} = -f$. \monades{9}
    \item Να αποδείξετε ότι $(g \circ f)(x) = \frac{1}{x}, x \in (0, +\infty)$. \monades{8}
    \item Έστω $h(x) = (g \circ f)(x)$. Να βρείτε τον μοναδικό αριθμό $\xi$ ο οποίος ικανοποιεί το συμπέρασμα του Θεωρήματος Μέσης Τιμής για την συνάρτηση $h$ στο διάστημα $[2, 8]$. \monades{8}
\end{alist}